\subsection{How a spacing is measured, and why the index needs a procedure}

The indices of Section~\ref{sec:problem} are only as good as the ruler that produces $z_i$. Three procedures appear in the work reviewed above, and they do not measure the same object.

The sticky-belt procedure lays a belt under the tube, moves the belt at the simulated ground speed, and reads gaps between seeds that have stuck where they fell \cite{panning2000,onal2012,onal2009}. It removes soil bounce and opener depth. It keeps tube bounce and meter errors. It is the right procedure for judging the feed-speed loop, because that loop's product is a release time. ISO~7256 is written in this spirit: a single-seed drill is tested for the regularity of the metered stream, not for the forest that may or may not grow \cite{iso7256}. The Chinese single-seed method GB/T~6973 organises the same three percentages, multiples, qualified, and misses, about the theoretical spacing \cite{gb6973}. JB/T~10293 then turns those percentages into a specification a manufacturer can claim against \cite{jb10293}. A hardwood machine should be willing to be scored by the same arithmetic even if it is not willing to adopt the maize pass marks. The arithmetic is what makes Table~\ref{tab:redmild} comparable with a future belt test.

The field procedure measures plants, not seeds. Emergence, rodents, and depth then enter $z_i$ \cite{lof2019,kachman1995}. A miss in the field is no longer evidence that the cell was empty. Kachman and Smith were explicit that a gap can be a failure to drop or a failure to emerge, and that the indices do not separate them. This manuscript therefore refuses field plant counts as the acceptance test of the controllers. They are the acceptance test of a plantation. Mixing the two is how a good meter is blamed for a dry week, and how a bad meter is excused by a kind season.

The optoelectronic procedure sits between the belt and the field. A gate in the tube times each seed and converts time to distance by multiplying by $v$ \cite{kocher1998,lan1999}. The conversion is exact only if $v$ is the speed used by the kinematic law and the seed's exit velocity along the row is negligible. A heavy acorn leaves the cell with a tangential speed $r\omega$. At $r=0.09$~m and $\omega=2.30$~rad/s that tangential speed is 0.21~m/s, which is not negligible beside a ground speed of 0.7~m/s. The horizontal component depends on the release angle. A tube gate that assumes the seed's row-speed equals $v$ will assign the wrong spacing by a few centimetres, which is enough to move a seed across the 0.5$X_t$ boundary when $X_t=0.25$~m. The estimation avoids that bias by placing seeds from the event model directly. A bench that uses a tube gate will need the release-angle correction derived in Section~\ref{sec:flight}, or it will disagree with the belt for a reason that is geometry rather than control.

Panning and colleagues also remind the reader that a laboratory ranking and a field ranking of planters need not match, even for sugarbeet, which is a far more uniform seed than an acorn \cite{panning2000}. The estimation in this manuscript is a third thing again: a ranking inside a declared model. It is closer to a belt test than to a field test, and it is not a belt test. The procedure that would make it a belt test is the one named at the end of Chapter~\ref{ch:estimate}.

\subsection{Particle forces underneath the factor map}

The logistic map is a summary. The physics it summarises is a contact problem. A discrete-element model represents each seed as a body with a position and a spin, and integrates the contact forces with its neighbours and with the wall \cite{cundall1979}. The normal force is a spring and a damper in the overlap. The tangential force is a spring limited by friction, which is what allows an arch to stand. For two acorns the overlap spring is a crude model of a wooden shell, but the structure of the argument does not depend on the spring law. An arch over the cell stands when the contact normals meet the wall at an angle steeper than the friction cone. Vibration that reverses the contact force for part of each cycle collapses the arch. Vibration that keeps the contacts sliding turns the hopper into a fluid and the cell fills with whatever arrives, including two small seeds.

Wu and colleagues used that style of model to show a filling angle that opens under vibration, and Hu and colleagues coupled it to an airflow solver for a pneumatic rapeseed meter \cite{wu2022,hu2021,lei2021}. The angle is a geometric observable of the same arch. A larger angle means contacts that have walked up the wall. The 10~Hz threshold of Zhang and colleagues is, in this language, the frequency at which the arch no longer reforms between cycles for a small seed \cite{zhang2024vib}. A large wooden seed has a longer contact time and a higher mass, so the frequency at which its arch collapses is lower. Equation~\eqref{eq:fstar} is that statement written as a decreasing function of diameter. It is not a modal analysis of an acorn. It is the smallest formula that gives the coordinator permission to treat pin oak and bur oak as different vibratory problems, which the single-frequency vegetable-seed result does not \cite{li2025drum}.

Inclination enters the same contact balance. On a slope the gravity vector gains a component along the throat, and the arch becomes asymmetric \cite{lei2022}. Wells could still meter red oak on a steep laboratory grade because the pneumatic force dominated that component for a seed that fitted the hole \cite{wells1990}. A seed that does not fit never receives the pneumatic force, and the slope component only makes the miss more likely. The slip term in Equation~\eqref{eq:zm} is the kinematic half of slope. The contact half is left inside the miss coefficient on $\delta_{+}$, which is why the estimation does not add a third slope mechanism. Adding one without a measurement would be a second unsupported constant.

This particle picture also explains a limit of vision. The camera sees the free surface and the seeds in the window. It does not see the contact forces under that surface. The bridge index is a proxy: high occupancy, low circularity, and oversize are the visible correlates of an arch. They will be wrong when an arch stands below a surface that looks open, which is a known hopper failure and a reason the constraint is capped. A capped constraint can be late. An uncapped constraint can stop the machine on a misread surface. The cap at 22 percent is the design's choice of which error is cheaper.
